
\section{The first nonaxisymmetric Fourier pair: pressure and harmonic leakage}
\label{sec:nonaxis-Fourier-XXIII}
This section studies a possible nonaxisymmetric escape from the
circulation obstruction. Fourier-mode selection, the cylindrical
vector Laplacian and Reynolds averaging are established prior art.
The source Navier--Stokes singularity is \emph{forced}; we do not
claim that a force-free construction has been completed.

Let $u_B=U(r,z,t)e_\theta$ be an axisymmetric pure-swirl background
(not assumed force-free). For
$v_m=e^{im\theta}(a e_r+b e_\theta+c e_z)$ put
$D_m=\partial_{rr}+r^{-1}\partial_r-m^2r^{-2}+\partial_{zz}$
and $L_m=\partial_t-\nu D_m$. Linearizing, with a
modewise right-hand side $H_m$, gives
\begin{align}
L_m a+\nu(a+2im b)/r^2+(imU/r)a-2Ub/r+\pi_{m,r}&=H_{m,r},\\
L_m b+\nu(b-2im a)/r^2+(imU/r)b+(U_r+U/r)a+U_zc
  +(im/r)\pi_m&=H_{m,\theta},\\
L_m c+(imU/r)c+\pi_{m,z}&=H_{m,z},\\
a_r+a/r+im b/r+c_z&=0.
\label{eq:nonaxis-system-XXIII}
\end{align}
The off-diagonal viscosity coefficients reflect the vector
Laplacian in its rotating basis and cannot be omitted.
Divergence and solenoidality give the necessary pressure equation
\begin{equation}\label{eq:nonaxis-pressure-XXIII}
D_m\pi_m=\operatorname{div}_mH_m
-\frac{2im}{r}(U_ra+U_zc)+\frac2r\partial_r(Ub),
\end{equation}
where
$\operatorname{div}_mH_m=H_{m,r,r}+r^{-1}H_{m,r}
+im r^{-1}H_{m,\theta}+H_{m,z,z}$.
This formula does not establish existence of an axis-regular
Poisson inverse with the original physical boundary data.

\begin{proposition}[A smooth first-harmonic source of mean angular torque]
\label{prop:nonaxis-torque-XXIII}
On $\mathbb R^2\times\mathbb S^1$ with axial period $2\pi$,
put $f(r)=e^{-r^2}$ and let $\lambda\in\mathbb R$.
For $m=+1$ set
\[
a=f\cos z,\qquad b=ia+\lambda r^2f\cos z,\qquad
c=-f'\sin z-i\lambda rf\sin z,
\]
and define the real velocity
$v=e^{i\theta}(a e_r+b e_\theta+c e_z)
+e^{-i\theta}(\bar a e_r+\bar b e_\theta+\bar c e_z)$.
It is smooth, finite-energy and divergence-free; in Cartesian form,
\[
v=(2f\cos z(1-\lambda xy),\,2\lambda x^2f\cos z,\,
       (4x+2\lambda y)f\sin z).
\]
Its azimuthal Reynolds tensor $Q=\langle v\otimes v\rangle_\theta$
has
\[
Q_{r\theta}=2\lambda r^2e^{-2r^2}\cos^2z,\qquad
Q_{z\theta}=\lambda r(2r^2-1)e^{-2r^2}\sin2z,
\]
and so the averaged angular convection satisfies
\begin{equation}\label{eq:nonaxis-az-force-XXIII}
g_\theta=\frac1{r^2}\partial_r(r^2Q_{r\theta})
             +\partial_zQ_{z\theta}
=2\lambda r e^{-2r^2}(2-2r^2+\cos2z).
\end{equation}
Its radial-slice torque is exactly
\begin{equation}\label{eq:nonaxis-slice-XXIII}
\int_0^\infty r^2g_\theta(r,z)\,dr
=\frac{\lambda}{4}\cos2z.
\end{equation}
At each point, $Q$ is positive semidefinite and has rank at most
two.
\end{proposition}
\begin{proof}
The modal divergence vanishes directly:
$a_r+a/r+ib/r+c_z=0$. Independently, the displayed Cartesian
field has zero divergence and is regular across the axis.
The angular average gives
$Q_{ij}=2\operatorname{Re}(v_{+,i}\overline{v_{+,j}})$,
a sum of two positive rank-one real tensors.
Substitution in the cylindrical tensor-divergence formula gives
\eqref{eq:nonaxis-az-force-XXIII}; exact Gaussian integration gives
\eqref{eq:nonaxis-slice-XXIII}. Direct real-angle integration
independently verifies both Reynolds tensor components.
\end{proof}

\begin{proposition}[Non-gradient second-harmonic residual]
\label{prop:nonaxis-mtwo-XXIII}
For the field in Proposition~\ref{prop:nonaxis-torque-XXIII}, the
$m=+2$ coefficient $N=(v_+\cdot\nabla)v_+$ of self-convection
has, at $\lambda=1$, $r=1$, $z=0$,
\begin{equation}\label{eq:nonaxis-mtwo-curl-XXIII}
(\nabla\times N)^{(+2)}_z=2(-2+5i)e^{-2}\ne0.
\end{equation}
Thus this \emph{particular} correction cannot be closed at
$m=0,\pm1$ by choosing scalar pressure.
\end{proposition}
\begin{proof}
In the rotating cylindrical basis, the $m=+2$ coefficients are
\[
N_r=aa_r+iba/r+ca_z-b^2/r,\quad
N_\theta=ab_r+ib^2/r+cb_z+ab/r,\quad
N_z=ac_r+ibc/r+cc_z.
\]
The axial curl coefficient is
$r^{-1}\partial_r(rN_\theta)-2iN_r/r$.
Exact differentiation gives \eqref{eq:nonaxis-mtwo-curl-XXIII}.
A scalar pressure gradient has zero curl.
\end{proof}

\paragraph{Connection to the existing H8 shell programme.}
The H8 inversion gate requires suitably small \emph{radial-slice}
torque in a slender shell, not merely zero total torque.
Equation~\eqref{eq:nonaxis-slice-XXIII} has zero average over $z$
but an $O(|\lambda|)$ nonzero slice torque. Any source-compatible
use of this mechanism requires $\lambda$ small at the right
slenderness scale or an independent cancelling profile.
The present manufactured domain is not the H8 physical source.

\paragraph{The full unforced and Lean gates.}
If the base residual lies in $m=0$, pure $m=\pm1$ perturbations
cannot cancel it linearly. Quadratic cross terms can generate
$m=0$ but the present witness generates $m=\pm2$ as well.
An exact zero-force theorem needs source-matched A.11/A.19
residuals and collars, a fully coupled pressure-compatible
harmonic inverse, nonlinear norm convergence and survival of
the terminal singularity. Independent SymPy and Cartesian
reconstructions, direct angular averages and 12,000 exact
rational finite tests pass. Lean establishes only finite
mode-index and Gram-minor algebra, not these PDE theorems.
